\documentclass[compress,10pt,xcolor=dvipsnames]{beamer}
\usepackage{pgfpages}
\setbeameroption{show notes on second screen}
\setbeamerfont{note page}{size=\footnotesize}

\usepackage[english]{babel}
\usepackage[utf8]{inputenc}
\usepackage{times}
\usepackage[T1]{fontenc}
\usepackage{alltt}

\usepackage{tikz}
\usetikzlibrary{calc,arrows,positioning,decorations.markings,shapes}
\usepackage{amsfonts}
\usepackage{amssymb}
\usepackage{amsmath}
\usepackage{graphicx}
\usepackage{tabularx}
\usepackage{hyperref}
\usepackage{natbib}
\setbeamertemplate{bibliography item}[text]

%% COLORS
\definecolor{awesome}{RGB}{3, 149, 222}
\definecolor{high1}{RGB}{3, 149, 222}
\definecolor{high2}{HTML}{F58137}
\definecolor{high3}{RGB}{120, 81, 169}
\definecolor{lightergrey}{rgb}{0.9, 0.9, 0.9}

%% TIKZ NODES & STYLES
\tikzset{
  boxnode/.style={draw=awesome!80, rectangle, rounded corners=2pt, minimum width=2.2cm, minimum height=1.0cm, align=center, thick},
  arrow/.style={>=stealth, thick, <->}
}

\usefonttheme{professionalfonts}
\usefonttheme{serif}

\mode<presentation>
{
    \usenavigationsymbolstemplate{}
    \setbeamertemplate{headline}
    {
        \begin{beamercolorbox}[sep=1ex,wd=\paperwidth]{whitebox}
            \hfill{\color{awesome}\insertframenumber}\hspace*{0.4em}
        \end{beamercolorbox}
    }

    \setbeamercovered{invisible}
    \setbeamercolor{frametitle}{fg=awesome}
    \setbeamercolor{date in head/foot}{fg=awesome}
    \setbeamercolor{itemize item}{fg=awesome}
    \setbeamercolor{itemize subitem}{fg=awesome}
    \setbeamercolor{itemize subsubitem}{fg=awesome}
    \setbeamercolor{enumerate item}{fg=awesome}

    \setbeamertemplate{itemize item}[square]
    \setbeamertemplate{itemize subitem}[circle]
    \setbeamertemplate{itemize subsubitem}[triangle]
}

\setbeamertemplate{footline}[text line]{%
  \parbox{\linewidth}{\vspace*{-6pt}{\color{lightergrey}\copyright Mohammad Abdulaziz}\hfill\insertshortauthor\hfill{\color{awesome}}}}
\setbeamertemplate{navigation symbols}{}

\newcommand{\chictitle}[2]{
  \begin{frame}[plain,noframenumbering]
    \begin{center}
      \vspace{2.5em}
      {\color{awesome}\Large\textbf{#1}} \\
      \vspace{0.8em}
      \small{#2}
      \vspace{3em}
      \begin{center}
        {\color{awesome}\large Mohammad Abdulaziz}
      \end{center}
    \end{center}
  \end{frame}
}


\setbeamersize{text margin left=8pt,text margin right=8pt}

\title[Proving Basic Algorithms Correct]{Proving Basic Algorithms Correct: Sorting & Locales}
\author[7CCSACTP]{Mohammad Abdulaziz}
\date{Week 6}

\begin{document}

\chictitle{Proving Basic Algorithms Correct: Sorting & Locales}{Computer-Assisted Theorem Proving (7CCSACTP)}

\begin{frame}{1. Formal Correctness of Merge Sort}
\note{
\begin{itemize}
  \item Specifying sortedness: \texttt{sorted} predicate
  \item Permutation & element preservation: \texttt{mset xs = mset ys}
  \item Proving correctness of list merging
\end{itemize}
}

\begin{itemize}
  \item \textbf{Specifying sortedness:}
    \begin{itemize}
      \item \texttt{sorted} predicate
    \end{itemize}
  \item \textbf{Permutation & element preservation:}
    \begin{itemize}
      \item \texttt{mset xs = mset ys}
    \end{itemize}
  \item \textbf{Proving correctness of list merging:} Proving correctness of list merging
\end{itemize}
\end{frame}

\begin{frame}{2. Complexity & Step Bounds}
\note{
\begin{itemize}
  \item Defining step-counting functions in Isabelle/HOL
  \item Setting up recurrence relations for recursive algorithms
  \item Proving $O(n \log n)$ step bounds via master theorem techniques
\end{itemize}
}

\begin{itemize}
  \item \textbf{Defining step-counting functions in Isabelle/HOL:} Defining step-counting functions in Isabelle/HOL
  \item \textbf{Setting up recurrence relations for recursive algorithms:} Setting up recurrence relations for recursive algorithms
  \item \textbf{Proving $O(n \log n)$ step bounds via master theorem techniques:} Proving $O(n \log n)$ step bounds via master theorem techniques
\end{itemize}
\end{frame}

\begin{frame}{3. Abstract Data Types via Locales}
\note{
\begin{itemize}
  \item Locales: structuring proof contexts and modular specifications
  \item Specifying abstract Set and Map ADTs
  \item Instantiating locales with concrete implementations
\end{itemize}
}

\begin{itemize}
  \item \textbf{Locales:}
    \begin{itemize}
      \item structuring proof contexts and modular specifications
    \end{itemize}
  \item \textbf{Specifying abstract Set and Map ADTs:} Specifying abstract Set and Map ADTs
  \item \textbf{Instantiating locales with concrete implementations:} Instantiating locales with concrete implementations
\end{itemize}
\end{frame}


\end{document}
